% !TEX root = task-7-3.tex
% Editable task bank. Append new BankTask blocks at the end of the document.
\documentclass[12pt,a4paper]{article}
\usepackage{physics-textbook}
\usepackage{graphicx,needspace}
\usetikzlibrary{decorations.pathmorphing}
\newif\ifShowSolutions
\ifdefined\TutorEdition
  \ShowSolutionstrue
\else
  \ShowSolutionsfalse % Change to \ShowSolutionstrue to display solutions.
\fi
\newcounter{banktask}
\newsavebox{\BankTaskBox}
\newcommand{\BankTaskContent}[3]{%
  \noindent\textbf{\thebanktask.}\quad #2\par
  \ifShowSolutions
    \Needspace{50mm}% Keep the solution heading with its opening text or diagram.
    #3\par
    \ifdefined\TutorEdition
      \ifcsname TutorSource#1\endcsname
        \nopagebreak\smallskip{\small\csname TutorSource#1\endcsname\par}%
      \fi
    \fi
  \fi
}
\newcommand{\BankTask}[3]{%
  \refstepcounter{banktask}\par\medskip
  \sbox{\BankTaskBox}{\begin{minipage}{\linewidth}%
    \BankTaskContent{#1}{#2}{#3}\end{minipage}}%
  \ifdim\dimexpr\ht\BankTaskBox+\dp\BankTaskBox\relax>\dimexpr\textheight-20mm\relax
    \Needspace{35mm}\BankTaskContent{#1}{#2}{#3}%
  \else
    \Needspace{\dimexpr\ht\BankTaskBox+\dp\BankTaskBox\relax}%
    \noindent\usebox{\BankTaskBox}\par
  \fi
  \medskip
}
\begin{document}
\BookSubsection{7.3}{Постановка эксперимента}


% Source 60DF4D: https://ege.fipi.ru/bank/questions.php?search=1&proj=BA1F39653304A5B041B656915DC36B38&qid=60DF4D
\BankTask{60DF4D}{
Необходимо собрать экспериментальную установку, с помощью которой можно определить плотность алюминия. Для этого школьник взял стакан с водой и алюминиевый шарик с крючком.

Какое дополнительное оборудование необходимо использовать для проведения этого эксперимента?
\begin{enumerate}
\item[1)] электронные весы
\item[2)] мензурка
\item[3)] секундомер
\item[4)] термометр
\item[5)] пружина
\end{enumerate}
}{
\smallskip\noindent\textbf{Решение от GPT}\par
Электронные весы (1) позволяют измерить массу шарика \(m\), а мензурка (2) — его объём по изменению уровня воды: \(V=V_2-V_1\). Плотность алюминия равна \(\rho=m/(V_2-V_1)\). Секундомер (3) и термометр (4) не нужны, а пружина (5) без известной жёсткости не позволяет измерить массу и не заменяет мензурку. Ответ: 1, 2.
}

% Source F68A6A: https://ege.fipi.ru/bank/questions.php?search=1&proj=BA1F39653304A5B041B656915DC36B38&qid=F68A6A
\BankTask{F68A6A}{
Необходимо собрать экспериментальную установку, с помощью которой можно определить плотность бензина. Для этого школьник взял стакан с бензином и динамометр.

Какие два предмета из приведённого ниже перечня оборудования необходимо дополнительно использовать для проведения этого эксперимента?
\begin{enumerate}
\item[1)] калориметр
\item[2)] мензурка
\item[3)] стальной цилиндр с крючком
\item[4)] термометр
\item[5)] пружина
\end{enumerate}
}{
\smallskip\noindent\textbf{Решение от GPT}\par
Нужны мензурка (2) и стальной цилиндр с крючком (3): объём цилиндра \(V\) определяют по объёму вытесненного бензина. Измерив динамометром вес цилиндра в воздухе \(P_1\) и при полном погружении в бензин \(P_2\), находят \(\rho=(P_1-P_2)/(gV)\). Калориметр (1) и термометр (4) не нужны, а пружина (5) не даёт необходимых измерений объёма погружённого тела. Ответ: 2, 3.
}

% Source 46D7FD: https://ege.fipi.ru/bank/questions.php?search=1&proj=BA1F39653304A5B041B656915DC36B38&qid=46D7FD
\BankTask{46D7FD}{
Необходимо собрать экспериментальную установку, с помощью которой можно определить жёсткость резинового жгута. Для этого школьник взял штатив с лапкой и резиновый жгут. Какие два предмета из приведённого ниже перечня оборудования необходимо дополнительно использовать для проведения этого эксперимента?
\begin{enumerate}
\item[1)] мензурка
\item[2)] линейка
\item[3)] секундомер
\item[4)] стакан с водой
\item[5)] динамометр
\end{enumerate}
}{
\smallskip\noindent\textbf{Решение от GPT}\par
Линейкой (2) измеряют удлинение жгута \(\Delta l=l-l_0\), а динамометром (5) — растягивающую силу \(F\). В области выполнения закона Гука жёсткость равна \(k=F/\Delta l\). Мензурка (1), секундомер (3) и стакан с водой (4) для этих измерений не нужны. Ответ: 2, 5.
}

% Source 9F7170: https://ege.fipi.ru/bank/questions.php?search=1&proj=BA1F39653304A5B041B656915DC36B38&qid=9F7170
\BankTask{9F7170}{
Необходимо собрать экспериментальную установку, с помощью которой можно определить коэффициент трения скольжения стали по дереву. Для этого школьник взял стальной брусок с крючком. Какие два предмета из приведённого ниже перечня оборудования необходимо дополнительно использовать для проведения этого эксперимента?
\begin{enumerate}
\item деревянная рейка
\item динамометр
\item мензурка
\item пластмассовая рейка
\item линейка
\end{enumerate}
}{
\smallskip\noindent\textbf{Решение от GPT}\par
Деревянная рейка (1) обеспечивает нужную поверхность, а динамометр (2) позволяет измерить вес бруска \(P\) и силу трения \(F_{\text{тр}}\) при равномерном движении по горизонтальной рейке. Тогда \(N=P\), и коэффициент трения равен \(\mu=F_{\text{тр}}/P\). Мензурка (3) и линейка (5) не нужны, а пластмассовая рейка (4) не подходит для исследования трения стали по дереву. Ответ: 12.
}

% Source 2C5D23: https://ege.fipi.ru/bank/questions.php?search=1&proj=BA1F39653304A5B041B656915DC36B38&qid=2C5D23
\BankTask{2C5D23}{
При помощи нитяного маятника необходимо экспериментально определить ускорение свободного падения. Для этого школьник взял штатив с муфтой и лапкой, нить и стальной шарик. Какие два предмета из приведённого ниже перечня оборудования необходимо дополнительно использовать для проведения этого эксперимента?
\begin{enumerate}
\item секундомер
\item линейка
\item мензурка
\item электронные весы
\item динамометр
\end{enumerate}
}{
\smallskip\noindent\textbf{Решение от GPT}\par
При малых колебаниях \(T=2\pi\sqrt{l/g}\), откуда \(g=4\pi^2l/T^2\). Секундомером (1) измеряют время \(t\) нескольких колебаний и находят \(T=t/n\), а линейкой (2) измеряют длину маятника от точки подвеса до центра шарика. Мензурка (3), электронные весы (4) и динамометр (5) не нужны: объём, масса и вес груза в эту формулу не входят. Ответ: 12.
}

% Source 1C49C9: https://ege.fipi.ru/bank/questions.php?search=1&proj=BA1F39653304A5B041B656915DC36B38&qid=1C49C9
\BankTask{1C49C9}{
Ученику необходимо на опыте обнаружить зависимость частоты свободных колебаний нитяного маятника от длины нити. У него имеется пять маятников, характеристики которых приведены в таблице. Какие два маятника необходимо взять ученику, чтобы провести данное исследование?
\begin{center}
\begin{tabular}{|c|c|c|c|}
\hline
\shortstack{№ маятника} & \shortstack{Длина нити,\\\(\text{см}\)} & \shortstack{Объём груза,\\\(\text{см}^{3}\)} & \shortstack{Материал, из которого\\сделан груз} \\
\hline
\(1\) & \(200\) & \(30\) & алюминий \\
\hline
\(2\) & \(150\) & \(60\) & алюминий \\
\hline
\(3\) & \(100\) & \(30\) & медь \\
\hline
\(4\) & \(100\) & \(60\) & медь \\
\hline
\(5\) & \(200\) & \(60\) & алюминий \\
\hline
\end{tabular}
\end{center}
}{
\smallskip\noindent\textbf{Решение от GPT}\par
Чтобы исследовать зависимость \(\nu=\frac{1}{2\pi}\sqrt{g/l}\) от длины нити, выбирают маятники с одинаковыми грузами и разными длинами нитей. У маятников 2 и 5 грузы алюминиевые, объёмом \(60\,\text{см}^3\), а длины нитей равны соответственно \(150\) и \(200\,\text{см}\). Ответ: 25.
}

% Source FA9CF6: https://ege.fipi.ru/bank/questions.php?search=1&proj=BA1F39653304A5B041B656915DC36B38&qid=FA9CF6
\BankTask{FA9CF6}{
Ученику необходимо на опыте обнаружить зависимость периода свободных колебаний пружинного маятника от массы груза. У него имеется пять пружинных маятников, характеристики которых приведены в таблице. Какие два маятника необходимо взять ученику, чтобы провести данное исследование? Грузы сплошные.
\begin{center}
\begin{tabular}{|c|c|c|c|}
\hline
\shortstack{№\\маятника} & \shortstack{Жёсткость пружины,\\\(\text{Н}/\text{м}\)} & \shortstack{Объём груза,\\$\text{см}^{3}$} & \shortstack{Материал, из которого\\сделан груз} \\
\hline
\(1\) & \(40\) & \(30\) & алюминий \\
\hline
\(2\) & \(60\) & \(60\) & алюминий \\
\hline
\(3\) & \(40\) & \(30\) & медь \\
\hline
\(4\) & \(10\) & \(30\) & алюминий \\
\hline
\(5\) & \(10\) & \(60\) & медь \\
\hline
\end{tabular}
\end{center}
}{
\smallskip\noindent\textbf{Решение от GPT}\par
Период пружинного маятника равен $T=2\pi\sqrt{m/k}$, поэтому нужно сохранить жёсткость пружины и изменить массу груза. У маятников 1 и 3 одинаковая жёсткость $k=40\,\text{Н/м}$, но разные массы $m=\rho V$, поскольку при одинаковом объёме материалы различаются. Ответ: 1 и 3. По приведённым данным также подходит пара 4 и 5: жёсткости одинаковы, а массы различны, поэтому выбор пары в условии не единственен.
}

% Source 0B0CF0: https://ege.fipi.ru/bank/questions.php?search=1&proj=BA1F39653304A5B041B656915DC36B38&qid=0B0CF0
\BankTask{0B0CF0}{
Необходимо экспериментально изучить зависимость ускорения тела, скользящего по шероховатой наклонной плоскости, от коэффициента трения груза о плоскость. На всех приведённых ниже рисунках указаны массы тел, углы наклона плоскостей к горизонту, коэффициенты трения. Какие две установки следует использовать для проведения такого исследования?
\begin{center}
\begin{tabular}{cc}
1) &
\begin{tikzpicture}[book diagram,scale=0.85]
\draw (0,0)--(0,2.08)--(3.6,0)--cycle;
\begin{scope}[shift={(1.3,1.329)},rotate=-30]
\draw (0,0) rectangle (0.85,0.48);
\node at (0.425,0.72) {$m$};
\end{scope}
\draw (2.85,0) arc[start angle=180,end angle=150,radius=0.75];
\node at (2.52,0.25) {$2\alpha$};
\node at (0.16,-0.32) {$2\mu$};
\draw[-{book tip}] (0.66,-0.27)--(1.61,1.15);
\end{tikzpicture}
\\[0.7em]
2) &
\begin{tikzpicture}[book diagram,scale=0.85]
\draw (0,0)--(0,0.965)--(3.6,0)--cycle;
\begin{scope}[shift={(1.25,0.630)},rotate=-15]
\draw (0,0) rectangle (0.85,0.48);
\node at (0.425,0.72) {$2m$};
\end{scope}
\draw (2.75,0) arc[start angle=180,end angle=165,radius=0.85];
\node at (2.56,0.16) {$\alpha$};
\node at (0.60,-0.32) {$\mu$};
\draw[-{book tip}] (0.92,-0.25)--(1.59,0.54);
\end{tikzpicture}
\\[0.7em]
3) &
\begin{tikzpicture}[book diagram,scale=0.85]
\draw (0,0)--(0,2.08)--(3.6,0)--cycle;
\begin{scope}[shift={(1.3,1.329)},rotate=-30]
\draw (0,0) rectangle (0.85,0.48);
\node at (0.425,0.72) {$3m$};
\end{scope}
\draw (2.85,0) arc[start angle=180,end angle=150,radius=0.75];
\node at (2.52,0.25) {$2\alpha$};
\node at (0.16,-0.32) {$2\mu$};
\draw[-{book tip}] (0.66,-0.27)--(1.61,1.15);
\end{tikzpicture}
\\[0.7em]
4) &
\begin{tikzpicture}[book diagram,scale=0.85]
\draw (0,0)--(0,0.965)--(3.6,0)--cycle;
\begin{scope}[shift={(1.25,0.630)},rotate=-15]
\draw (0,0) rectangle (0.85,0.48);
\node at (0.425,0.72) {$m$};
\end{scope}
\draw (2.75,0) arc[start angle=180,end angle=165,radius=0.85];
\node at (2.56,0.16) {$\alpha$};
\node at (0.60,-0.32) {$2\mu$};
\draw[-{book tip}] (0.96,-0.25)--(1.59,0.54);
\end{tikzpicture}
\\[0.7em]
5) &
\begin{tikzpicture}[book diagram,scale=0.85]
\draw (0,0)--(0,0.965)--(3.6,0)--cycle;
\begin{scope}[shift={(1.25,0.630)},rotate=-15]
\draw (0,0) rectangle (0.85,0.48);
\node at (0.425,0.72) {$m$};
\end{scope}
\draw (2.75,0) arc[start angle=180,end angle=165,radius=0.85];
\node at (2.56,0.16) {$\alpha$};
\node at (0.60,-0.32) {$\mu$};
\draw[-{book tip}] (0.92,-0.25)--(1.59,0.54);
\end{tikzpicture}
\end{tabular}
\end{center}
}{
\smallskip\noindent\textbf{Решение от GPT}\par
Для скольжения вниз $a=g(\sin\alpha-\mu\cos\alpha)$; в опытах следует изменять коэффициент трения, сохраняя массу и угол наклона. Установки 4 и 5 имеют одинаковые массу $m$ и угол $\alpha$, но коэффициенты трения равны $2\mu$ и $\mu$. Установки 1 и 3 имеют угол $2\alpha$ и одинаковый коэффициент трения, а установка 2 отличается от установок 4 и 5 массой $2m$. Ответ: 4 и 5.
}

% Source A6D7D1: https://ege.fipi.ru/bank/questions.php?search=1&proj=BA1F39653304A5B041B656915DC36B38&qid=A6D7D1
\BankTask{A6D7D1}{
Ученик изучает силу Архимеда, действующую на тела, полностью погружённые в жидкость. В его распоряжении имеется пять установок, состоящих из ёмкости с жидкостью и сплошного шарика. Какие две из перечисленных в таблице установок необходимы ученику для того, чтобы на опыте обнаружить, зависит ли сила Архимеда от плотности материала шарика?
\begin{center}
\begin{tabular}{|c|c|c|c|}
\hline
\shortstack{№\\установки} & \shortstack{Жидкость, налитая\\в ёмкость} & Объём шарика & \shortstack{Материал, из которого\\сделан шарик} \\
\hline
1 & вода & $30\,\text{см}^{3}$ & сталь \\
\hline
2 & вода & $20\,\text{см}^{3}$ & алюминий \\
\hline
3 & керосин & $20\,\text{см}^{3}$ & алюминий \\
\hline
4 & подсолнечное масло & $30\,\text{см}^{3}$ & сталь \\
\hline
5 & вода & $30\,\text{см}^{3}$ & алюминий \\
\hline
\end{tabular}
\end{center}
}{
\smallskip\noindent\textbf{Решение от GPT}\par
Нужно взять шарики одинакового объёма из разных материалов и полностью погрузить их в одну и ту же жидкость. Этим требованиям соответствуют установки 1 и 5: вода, объём $30\,\text{см}^{3}$, материалы — сталь и алюминий. Сила Архимеда $F_A=\rho_{\text{ж}}gV$ в обоих опытах одинакова и не зависит от плотности материала шарика. Ответ: 1 и 5.
}

% Source 2493F5: https://ege.fipi.ru/bank/questions.php?search=1&proj=BA1F39653304A5B041B656915DC36B38&qid=2493F5
\BankTask{2493F5}{
Необходимо собрать экспериментальную установку, с помощью которой можно определить оптическую силу собирающей линзы. В качестве источника света школьник взял горящую свечу. Какие два предмета из приведённого ниже перечня оборудования необходимо дополнительно использовать для проведения этого эксперимента?
\begin{enumerate}
\item гальванометр
\item линейка
\item экран
\item зеркало
\item динамометр
\end{enumerate}
}{
\smallskip\noindent\textbf{Решение от GPT}\par
На экране (3) получают резкое изображение свечи, а линейкой (2) измеряют расстояния от свечи до линзы \(a\) и от линзы до экрана \(b\). Оптическая сила линзы \(D=1/F=1/a+1/b\); при расстояниях в метрах она выражается в диоптриях. Гальванометр (1) измеряет силу тока, зеркало (4) для этого опыта не требуется, динамометр (5) измеряет силу. Ответ: 2, 3.
}

% Source 154E79: https://ege.fipi.ru/bank/questions.php?search=1&proj=BA1F39653304A5B041B656915DC36B38&qid=154E79
\BankTask{154E79}{
Необходимо собрать экспериментальную установку, с помощью которой можно определить показатель преломления стекла. Для этого школьник взял источник света, создающий узкий пучок света, карандаш и циркуль. Какие два предмета из приведённого ниже перечня оборудования необходимо дополнительно использовать для проведения этого эксперимента?
\begin{enumerate}
\item плоскопараллельная стеклянная пластина
\item линейка
\item плоскопараллельная плексигласовая пластина
\item зеркало
\item собирающая линза
\end{enumerate}
}{
\smallskip\noindent\textbf{Решение от GPT}\par
Нужны стеклянная пластина (1) и линейка (2): отмечают ход луча, строят нормаль к границе и с помощью циркуля и линейки определяют отношение \(n=\sin\alpha/\sin\beta\), считая показатель преломления воздуха равным единице. Плексигласовая пластина (3) позволяет исследовать другой материал, зеркало (4) даёт отражение, а собирающая линза (5) не нужна для измерения углов на плоской границе стекла. Ответ: 1, 2.
}

% Source c70327: https://ege.fipi.ru/bank/questions.php?search=1&proj=BA1F39653304A5B041B656915DC36B38&qid=c70327
\BankTask{c70327}{
Необходимо на опыте обнаружить зависимость объёма газа, находящегося в сосуде под подвижным поршнем, от температуры газа. Имеются пять различных сосудов с манометрами и термометрами. Сосуды наполнены одним и тем же газом при различных температурах и давлениях (см. таблицу). Какие два сосуда необходимо взять, чтобы провести это исследование?
\begin{center}
\begin{tabular}{|c|c|c|c|}
\hline
\shortstack{№\\сосуда} & \shortstack{Давление\\газа, \(\text{кПа}\)} & \shortstack{Температура газа\\в сосуде, \({}^{\circ}\mathrm{C}\)} & Масса газа, \(\text{г}\) \\
\hline
\(1\) & \(80\) & \(50\) & \(15\) \\
\hline
\(2\) & \(100\) & \(80\) & \(10\) \\
\hline
\(3\) & \(120\) & \(100\) & \(10\) \\
\hline
\(4\) & \(80\) & \(50\) & \(20\) \\
\hline
\(5\) & \(100\) & \(120\) & \(10\) \\
\hline
\end{tabular}
\end{center}
}{
\smallskip\noindent\textbf{Решение от GPT}\par
Для исследования зависимости объёма от температуры нужно сохранить одинаковыми массу газа и давление, изменив температуру. Из уравнения \(pV=(m/M)RT\) при этих условиях следует \(V/T=\mathrm{const}\). Подходят сосуды 2 и 5: в обоих давление \(100\,\text{кПа}\) и масса газа \(10\,\text{г}\), а температуры различаются. Ответ: 2, 5.
}

% Source cBe7DB: https://ege.fipi.ru/bank/questions.php?search=1&proj=BA1F39653304A5B041B656915DC36B38&qid=cBe7DB
\BankTask{cBe7DB}{
Ученику необходимо на опыте обнаружить зависимость объёма газа, находящегося в сосуде под подвижным поршнем, от молярной массы газа.

У него имеется пять различных сосудов с манометрами и термометрами. Сосуды наполнены равными массами различных газов при различных температурах (см. таблицу). Какие два сосуда необходимо взять ученику, чтобы провести исследование?

\begin{center}
\begin{tabular}{|c|c|c|c|}
\hline
\shortstack{№\\сосуда} & Давление, \(\text{кПа}\) & \shortstack{Температура газа\\в сосуде, \(\text{К}\)} & Газ в сосуде \\
\hline
\(1\) & \(100\) & \(280\) & кислород \\
\hline
\(2\) & \(50\) & \(270\) & кислород \\
\hline
\(3\) & \(100\) & \(280\) & азот \\
\hline
\(4\) & \(50\) & \(300\) & азот \\
\hline
\(5\) & \(60\) & \(320\) & кислород \\
\hline
\end{tabular}
\end{center}
}{
\smallskip\noindent\textbf{Решение от GPT}\par
Из уравнения состояния идеального газа \(V=\dfrac{mRT}{Mp}\): для исследования зависимости объёма от молярной массы \(M\) нужно сохранить одинаковыми массу, температуру и давление газа. Массы всех газов равны, а в сосудах \(1\) и \(3\) одинаковы давление \(100\,\text{кПа}\) и температура \(280\,\text{К}\), но различны газы и их молярные массы. Ответ: \(13\).
}

% Source 328FBF: https://ege.fipi.ru/bank/questions.php?search=1&proj=BA1F39653304A5B041B656915DC36B38&qid=328FBF
\BankTask{328FBF}{
Для лабораторной работы по обнаружению зависимости сопротивления проводника от его длины ученику выдали пять проводников, характеристики которых указаны в таблице. Какие два из предложенных ниже проводников необходимо взять ученику, чтобы провести данное исследование?

\begin{center}
\begin{tabular}{|c|c|c|c|}
\hline
№ проводника & \shortstack{Длина проводника,\\\(\text{см}\)} & \shortstack{Диаметр проводника,\\\(\text{мм}\)} & Материал \\
\hline
\(1\) & \(200\) & \(1{,}0\) & алюминий \\
\hline
\(2\) & \(100\) & \(0{,}5\) & медь \\
\hline
\(3\) & \(100\) & \(1{,}0\) & медь \\
\hline
\(4\) & \(100\) & \(0{,}5\) & алюминий \\
\hline
\(5\) & \(200\) & \(1{,}0\) & медь \\
\hline
\end{tabular}
\end{center}
}{
\smallskip\noindent\textbf{Решение от GPT}\par
Сопротивление проводника \(R=\rho l/S\), поэтому при изменении длины нужно сохранить одинаковыми материал и диаметр проводников. Проводники \(3\) и \(5\) изготовлены из меди, имеют диаметр \(1{,}0\,\text{мм}\) и разные длины: \(100\,\text{см}\) и \(200\,\text{см}\). Ответ: \(35\).
}

% Source 837BC0: https://ege.fipi.ru/bank/questions.php?search=1&proj=BA1F39653304A5B041B656915DC36B38&qid=837BC0
\BankTask{837BC0}{
Для проведения лабораторной работы по обнаружению зависимости сопротивления проводника от его диаметра ученику выдали пять проводников, характеристики которых приведены в таблице. Какие два из предложенных ниже проводников необходимо взять ученику, чтобы провести данное исследование?

\begin{center}
\begin{tabular}{|c|c|c|c|}
\hline
№ проводника & Длина проводника & Диаметр проводника & Материал \\
\hline
\(1\) & \(100\,\text{см}\) & \(1{,}0\,\text{мм}\) & медь \\
\hline
\(2\) & \(200\,\text{см}\) & \(0{,}5\,\text{мм}\) & медь \\
\hline
\(3\) & \(200\,\text{см}\) & \(1{,}0\,\text{мм}\) & медь \\
\hline
\(4\) & \(100\,\text{см}\) & \(0{,}5\,\text{мм}\) & алюминий \\
\hline
\(5\) & \(300\,\text{см}\) & \(1{,}0\,\text{мм}\) & медь \\
\hline
\end{tabular}
\end{center}
}{
\smallskip\noindent\textbf{Решение от GPT}\par
Сопротивление проводника \(R=\dfrac{4\rho l}{\pi d^2}\), поэтому для исследования зависимости от диаметра нужно выбрать проводники одинаковых длины и материала. Проводники \(2\) и \(3\) изготовлены из меди и имеют длину \(200\,\text{см}\), а их диаметры различны: \(0{,}5\,\text{мм}\) и \(1{,}0\,\text{мм}\). Ответ: \(23\).
}

% Source D860B0: https://ege.fipi.ru/bank/questions.php?search=1&proj=BA1F39653304A5B041B656915DC36B38&qid=D860B0
\BankTask{D860B0}{
Ученик изучает законы постоянного тока. В его распоряжении имеется пять аналогичных электрических цепей (см. рисунок) с различными источниками и внешними сопротивлениями, характеристики которых указаны в таблице. Какие две цепи необходимо взять ученику для того, чтобы на опыте обнаружить зависимость силы тока, протекающего в цепи, от внутреннего сопротивления источника?

\begin{center}
\begin{tabular}{|c|c|c|c|}
\hline
№ цепи & \shortstack{ЭДС источника\\$\varepsilon$, \(\text{В}\)} & \shortstack{Внутреннее сопротивление\\источника $r$, \(\text{Ом}\)} & \shortstack{Внешнее сопротивление\\$R$, \(\text{Ом}\)} \\
\hline
\(1\) & \(6\) & \(1{,}0\) & \(5\) \\
\hline
\(2\) & \(9\) & \(0{,}5\) & \(8\) \\
\hline
\(3\) & \(8\) & \(1{,}5\) & \(16\) \\
\hline
\(4\) & \(9\) & \(1{,}5\) & \(8\) \\
\hline
\(5\) & \(6\) & \(0{,}5\) & \(16\) \\
\hline
\end{tabular}
\end{center}

\begin{center}
\begin{tikzpicture}[book diagram]
\draw (0,0.82) -- (0,1.4) -- (0.9,1.4);
\draw (0,0.38) -- (0,0.02) -- (0.5,0.02);
\draw (0,0.6) circle (0.22);
\node at (0,0.6) {$\mathrm{A}$};
\draw (0.9,1.08) -- (0.9,1.72);
\draw (1.06,1.25) -- (1.06,1.55);
\draw (1.06,1.4) -- (1.95,1.4) -- (1.95,0.02) -- (1.5,0.02);
\draw (0.5,-0.18) rectangle (1.5,0.22);
\node[above] at (1,0.22) {$R$};
\node[above] at (1.4,1.55) {$\mathcal{E},\,r$};
\end{tikzpicture}
\end{center}
}{
\smallskip\noindent\textbf{Решение от GPT}\par
По закону Ома для полной цепи $I=\dfrac{\varepsilon}{R+r}$. Для исследования зависимости от $r$ необходимо сохранить одинаковыми ЭДС и внешнее сопротивление: у цепей \(2\) и \(4\) $\varepsilon=9\,\text{В}$ и $R=8\,\text{Ом}$, а внутренние сопротивления различаются. Ответ: \(24\).
}

% Source A88c1A: https://ege.fipi.ru/bank/questions.php?search=1&proj=BA1F39653304A5B041B656915DC36B38&qid=A88c1A
\BankTask{A88c1A}{
Необходимо экспериментально обнаружить зависимость силы тока, протекающего в цепи, от ЭДС источника. Какие две электрические цепи, схемы которых приведены на рисунках, следует использовать для такого исследования?

\begin{center}
\begin{tikzpicture}[book diagram]
\begin{scope}
\node at (0.75,1.95) {1)};
\draw (0,0.77) -- (0,1.3) -- (0.72,1.3);
\draw (0,0.33) -- (0,0) -- (0.38,0);
\draw (0,0.55) circle (0.22);
\node at (0,0.55) {$\mathrm{A}$};
\draw (0.72,1.03) -- (0.72,1.57);
\draw (0.86,1.17) -- (0.86,1.43);
\draw (0.86,1.3) -- (1.6,1.3) -- (1.6,0) -- (1.22,0);
\draw (0.38,-0.18) rectangle (1.22,0.18);
\node[above] at (0.8,0.18) {$R$};
\node[above] at (1.17,1.43) {$\mathcal{E},r$};
\end{scope}
\begin{scope}[xshift=2.6cm]
\node at (0.75,1.95) {2)};
\draw (0,0.77) -- (0,1.3) -- (0.72,1.3);
\draw (0,0.33) -- (0,0) -- (0.38,0);
\draw (0,0.55) circle (0.22);
\node at (0,0.55) {$\mathrm{A}$};
\draw (0.72,1.03) -- (0.72,1.57);
\draw (0.86,1.17) -- (0.86,1.43);
\draw (0.86,1.3) -- (1.6,1.3) -- (1.6,0) -- (1.22,0);
\draw (0.38,-0.18) rectangle (1.22,0.18);
\node[above] at (0.8,0.18) {$2R$};
\node[above] at (1.17,1.43) {$2\mathcal{E},r$};
\end{scope}
\begin{scope}[xshift=5.2cm]
\node at (0.75,1.95) {3)};
\draw (0,0.77) -- (0,1.3) -- (0.72,1.3);
\draw (0,0.33) -- (0,0) -- (0.38,0);
\draw (0,0.55) circle (0.22);
\node at (0,0.55) {$\mathrm{A}$};
\draw (0.72,1.03) -- (0.72,1.57);
\draw (0.86,1.17) -- (0.86,1.43);
\draw (0.86,1.3) -- (1.6,1.3) -- (1.6,0) -- (1.22,0);
\draw (0.38,-0.18) rectangle (1.22,0.18);
\node[above] at (0.8,0.18) {$R$};
\node[above] at (1.17,1.43) {$\mathcal{E},2r$};
\end{scope}
\begin{scope}[xshift=7.8cm]
\node at (0.75,1.95) {4)};
\draw (0,0.77) -- (0,1.3) -- (0.72,1.3);
\draw (0,0.33) -- (0,0) -- (0.38,0);
\draw (0,0.55) circle (0.22);
\node at (0,0.55) {$\mathrm{A}$};
\draw (0.72,1.03) -- (0.72,1.57);
\draw (0.86,1.17) -- (0.86,1.43);
\draw (0.86,1.3) -- (1.6,1.3) -- (1.6,0) -- (1.22,0);
\draw (0.38,-0.18) rectangle (1.22,0.18);
\node[above] at (0.8,0.18) {$2R$};
\node[above] at (1.17,1.43) {$2\mathcal{E},2r$};
\end{scope}
\begin{scope}[xshift=10.4cm]
\node at (0.75,1.95) {5)};
\draw (0,0.77) -- (0,1.3) -- (0.72,1.3);
\draw (0,0.33) -- (0,0) -- (0.38,0);
\draw (0,0.55) circle (0.22);
\node at (0,0.55) {$\mathrm{A}$};
\draw (0.72,1.03) -- (0.72,1.57);
\draw (0.86,1.17) -- (0.86,1.43);
\draw (0.86,1.3) -- (1.6,1.3) -- (1.6,0) -- (1.22,0);
\draw (0.38,-0.18) rectangle (1.22,0.18);
\node[above] at (0.8,0.18) {$2R$};
\node[above] at (1.17,1.43) {$\mathcal{E},r$};
\end{scope}
\end{tikzpicture}
\end{center}
}{
\smallskip\noindent\textbf{Решение от GPT}\par
Согласно формуле $I=\dfrac{\mathcal{E}}{R_{\text{внеш}}+r_{\text{внутр}}}$, для проверки зависимости от ЭДС оба сопротивления должны быть одинаковыми. Цепи 2 и 5 имеют внешнее сопротивление $2R$ и внутреннее сопротивление $r$, а их ЭДС равны $2\mathcal{E}$ и $\mathcal{E}$ соответственно. Для цепи 1 нет пары с теми же сопротивлениями и другой ЭДС; цепи 3 и 4 также не имеют таких пар, поскольку их внешние сопротивления различаются. Ответ: 25.
}

% Source 82e66e: https://ege.fipi.ru/bank/questions.php?search=1&proj=BA1F39653304A5B041B656915DC36B38&qid=82e66e
\BankTask{82e66e}{
Ученику необходимо на опыте обнаружить зависимость электроёмкости плоского конденсатора от площади его пластин. У него имеется пять конденсаторов, характеристики которых приведены в таблице.

Какие два конденсатора необходимо взять ученику, чтобы провести данное исследование?

\begin{center}
\begin{tabular}{|c|c|c|c|}
\hline
\shortstack{№\\конденсатора} & \shortstack{Расстояние между\\пластинами, \(\text{мм}\)} & \shortstack{Площадь пластин,\\$\text{см}^{2}$} & \shortstack{Диэлектрик, заполняющий\\пространство между пластинами} \\
\hline
\(1\) & \(0{,}6\) & \(20\) & парафин \\
\hline
\(2\) & \(0{,}8\) & \(10\) & парафин \\
\hline
\(3\) & \(0{,}9\) & \(25\) & слюда \\
\hline
\(4\) & \(0{,}6\) & \(15\) & парафин \\
\hline
\(5\) & \(1{,}0\) & \(10\) & слюда \\
\hline
\end{tabular}
\end{center}
}{
\smallskip\noindent\textbf{Решение от GPT}\par
Электроёмкость плоского конденсатора $C=\varepsilon_0\varepsilon S/d$. Следует изменять площадь $S$, сохраняя расстояние $d$ и диэлектрик одинаковыми: этому условию соответствуют конденсаторы 1 и 4, у которых $d=0{,}6\,\text{мм}$, диэлектрик — парафин, а площади равны $20$ и $15\,\text{см}^2$. Ответ: 14.
}

% Source 7D4F20: https://ege.fipi.ru/bank/questions.php?search=1&proj=BA1F39653304A5B041B656915DC36B38&qid=7D4F20
\BankTask{7D4F20}{
Необходимо экспериментально изучить зависимость заряда, накопленного конденсатором, от сопротивления резистора. Какие две схемы следует использовать для проведения такого исследования?

\begin{center}
\begin{tikzpicture}[book diagram, x=1cm, y=1cm]
\foreach \x/\n/\source/\res in {0/1/{\mathcal{E},r}/{R},2.6/2/{2\mathcal{E},2r}/{2R},5.2/3/{2\mathcal{E},r}/{2R},7.8/4/{\mathcal{E},2r}/{R},10.4/5/{\mathcal{E},r}/{2R}} {
\begin{scope}[shift={(\x,0)}]
\node at (0.65,-0.55) {\n)};
\draw (0,0)--(0,1.8)--(0.57,1.8);
\draw (0.72,1.8)--(1.35,1.8)--(1.35,0)--(0.73,0);
\draw (0.57,1.57)--(0.57,2.03);
\draw (0.72,1.68)--(0.72,1.92);
\node[anchor=south west] at (0.67,1.97) {$\source$};
\draw (0,0.83)--(0.3,0.83);
\draw (0.3,0.68) rectangle (1.04,0.98);
\draw (1.04,0.83)--(1.35,0.83);
\fill (0,0.83) circle (0.035);
\fill (1.35,0.83) circle (0.035);
\node[anchor=south] at (0.67,1.01) {$\res$};
\draw (0,0)--(0.57,0);
\draw (0.57,-0.24)--(0.57,0.24);
\draw (0.73,-0.24)--(0.73,0.24);
\node[anchor=south west] at (0.76,0.04) {$C$};
\end{scope}
}
\end{tikzpicture}
\end{center}
}{
\smallskip\noindent\textbf{Решение от GPT}\par
Менять нужно только сопротивление резистора, сохраняя одинаковыми ЭДС источника, его внутреннее сопротивление и ёмкость конденсатора. Этому условию соответствуют схемы 1 и 5: сопротивления равны соответственно $R$ и $2R$, а остальные параметры совпадают. В схеме 2 изменены ЭДС и внутреннее сопротивление, в схеме 3 — ЭДС, в схеме 4 — внутреннее сопротивление относительно схемы 1. После установления режима $q=CU=\dfrac{C\mathcal{E}R_{\text{рез}}}{R_{\text{рез}}+r}$. Ответ: 15.
}

% Source 70B4D7: https://ege.fipi.ru/bank/questions.php?search=1&proj=BA1F39653304A5B041B656915DC36B38&qid=70B4D7
\BankTask{70B4D7}{
Ученику необходимо обнаружить зависимость периода свободных электромагнитных колебаний в колебательном контуре от индуктивности катушки. Какие два колебательных контура он должен выбрать для проведения такого опыта?

\begin{center}
\begin{tikzpicture}[book diagram, x=1cm, y=1cm]
\foreach \x/\n/\ind/\cap in {0/1/{2L}/{C},3/2/{L}/{C},6/3/{L}/{2C},9/4/{2L}/{3C},12/5/{\dfrac{L}{2}}/{\dfrac{C}{2}}} {
\begin{scope}[shift={(\x,0)}]
\node at (0.5,-0.45) {\n)};
\draw (0,0.28)--(0,0)--(1,0)--(1,0.66);
\draw (0,1.09)--(0,1.4)--(1,1.4)--(1,0.82);
\draw plot[domain=0:1800,samples=241,variable=\t] ({0.10*sin(\t)},{0.235+0.81*\t/1800+0.045*cos(\t)});
\draw (0.74,0.66)--(1.26,0.66);
\draw (0.74,0.82)--(1.26,0.82);
\node[anchor=east] at (-0.15,0.7) {$\ind$};
\node[anchor=west] at (1.3,0.74) {$\cap$};
\end{scope}
}
\end{tikzpicture}
\end{center}
}{
\smallskip\noindent\textbf{Решение от GPT}\par
По формуле Томсона $T=2\pi\sqrt{LC}$, поэтому для исследования зависимости периода от индуктивности необходимо изменять $L$ при постоянной ёмкости. В контурах 1 и 2 ёмкость одинакова и равна $C$, а индуктивности равны $2L$ и $L$; следовательно, $T_1/T_2=\sqrt{2}$. В контурах 3, 4 и 5 ёмкости равны соответственно $2C$, $3C$ и $C/2$, поэтому они не образуют подходящей пары с одинаковой ёмкостью. Ответ: 12.
}

\end{document}
